% Options for packages loaded elsewhere
\PassOptionsToPackage{unicode}{hyperref}
\PassOptionsToPackage{hyphens}{url}
%
\documentclass[
  11pt,
]{article}
\usepackage{lmodern}
\usepackage{amssymb,amsmath}
\usepackage{ifxetex,ifluatex}
\ifnum 0\ifxetex 1\fi\ifluatex 1\fi=0 % if pdftex
  \usepackage[T1]{fontenc}
  \usepackage[utf8]{inputenc}
  \usepackage{textcomp} % provide euro and other symbols
\else % if luatex or xetex
  \usepackage{unicode-math}
  \defaultfontfeatures{Scale=MatchLowercase}
  \defaultfontfeatures[\rmfamily]{Ligatures=TeX,Scale=1}
\fi
% Use upquote if available, for straight quotes in verbatim environments
\IfFileExists{upquote.sty}{\usepackage{upquote}}{}
\IfFileExists{microtype.sty}{% use microtype if available
  \usepackage[]{microtype}
  \UseMicrotypeSet[protrusion]{basicmath} % disable protrusion for tt fonts
}{}
\makeatletter
\@ifundefined{KOMAClassName}{% if non-KOMA class
  \IfFileExists{parskip.sty}{%
    \usepackage{parskip}
  }{% else
    \setlength{\parindent}{0pt}
    \setlength{\parskip}{6pt plus 2pt minus 1pt}}
}{% if KOMA class
  \KOMAoptions{parskip=half}}
\makeatother
\usepackage{xcolor}
\IfFileExists{xurl.sty}{\usepackage{xurl}}{} % add URL line breaks if available
\IfFileExists{bookmark.sty}{\usepackage{bookmark}}{\usepackage{hyperref}}
\hypersetup{
  hidelinks,
  pdfcreator={LaTeX via pandoc}}
\urlstyle{same} % disable monospaced font for URLs
\usepackage[margin=1in]{geometry}
\setlength{\emergencystretch}{3em} % prevent overfull lines
\providecommand{\tightlist}{%
  \setlength{\itemsep}{0pt}\setlength{\parskip}{0pt}}
\setcounter{secnumdepth}{-\maxdimen} % remove section numbering
\usepackage{esint}
\setmathfont{Asana Math}
\setlength{\emergencystretch}{3em}

\author{}
\date{}

\begin{document}

\hypertarget{half-flat-bubble-exclusion-beyond-aspherical-topology}{%
\section{Half-flat bubble exclusion beyond aspherical
topology}\label{half-flat-bubble-exclusion-beyond-aspherical-topology}}

23 September 2026. This continuation weakens the topology in the
compact-core argument, and therefore in the compact-regime realization
theorem. It does not alter the archived proofs. The classification of
half-flat ALE spaces is a classical input; the extension below
identifies the homotopical information actually needed for their
exclusion.

\hypertarget{the-topological-hypothesis}{%
\subsection{1. The topological
hypothesis}\label{the-topological-hypothesis}}

Let \(M\) be a closed connected smooth oriented four-manifold satisfying
\[
 \pi_1(M)\text{ is torsion-free},\qquad \pi_2(M)=0.       \tag{H}
\] Equivalently, the second condition says
\(H_2(\widetilde M;\mathbb Z)=0\): the universal cover is simply
connected, so this is the degree-two Hurewicz theorem. No vanishing of
\(\pi_3(M)\), no contractibility, and no assumption on spin are
required.

\textbf{Theorem 1 (compact-core exclusion under (H)).} No sufficiently
large compact core of a nonflat smooth Ricci-flat ALE four-manifold with
one identically vanishing Weyl half embeds smoothly in a manifold
satisfying (H). The assertion is independent of the orientation of the
proposed embedding.

\textbf{Proof.} Reverse orientation on the ALE space if necessary, so
that \(W^+=0\). Let \(B\) be a large core with one connected spherical
space-form boundary \(Y\), and with the homotopy type of the ALE space.
The classification of Wright, together with Suvaina's description of the
Kähler ALE quotients, gives the following alternatives:

\begin{enumerate}
\def\labelenumi{\arabic{enumi}.}
\tightlist
\item
  \(\pi_1(B)\) is finite and the rational intersection form of \(B\) is
  negative definite, with \(b_2(B)>0\); or
\item
  \(B\) is a rational homology ball with \[
  \pi_1(B)=H_1(B;\mathbb Z)=\mathbb Z_n,
  \qquad Y=L(n^2,nr-1),\quad n\ge2,\ (n,r)=1.
  \tag{1}
  \]
\end{enumerate}

These are the same classification inputs used in the
\href{../mathematics/current/line1_aspherical_no_bubbles.md}{inherited
compact-core proof}. The nonflat assumption excludes the Euclidean case.

Suppose \(i:B\hookrightarrow M\) is an embedding. The image of the
finite group \(\pi_1(B)\) in the torsion-free group \(\pi_1(M)\) is
trivial. Thus \(i\) lifts to an embedding into \(\widetilde M\). The
lift is injective since its composition with the covering projection is
injective. This does not assert that a limiting fundamental group
injects into the fundamental group of \(M\).

In the first case a compact rational two-cycle in the interior of \(B\)
has nonzero self-intersection. Its image is zero in
\(H_2(\widetilde M;\mathbb Q)\), which vanishes by (H). Intersection of
compact cycles is unchanged by an embedding, up to an overall
orientation sign, and is zero for a null-homologous compact cycle. This
is a contradiction.

In the second case put \(C=\widetilde M\setminus\operatorname{int}B\)
and apply integral Mayer--Vietoris using collar enlargements. Since
\(H_1(\widetilde M;\mathbb Z)=H_2(\widetilde M;\mathbb Z)=0\), \[
 \mathbb Z_{n^2}\cong H_1(Y;\mathbb Z)
 \cong H_1(B;\mathbb Z)\oplus H_1(C;\mathbb Z)
 \cong\mathbb Z_n\oplus H_1(C;\mathbb Z).               \tag{2}
\] The last summand has order \(n\). The right side has exponent at most
\(n\), whereas the left side has exponent \(n^2\). This contradiction
also works for composite \(n\). Only the first two homology groups of
the universal cover were used. \(\square\)

\hypertarget{what-an-arbitrary-remaining-ale-core-would-have-to-be}{%
\subsubsection{1.1 What an arbitrary remaining ALE core would have to
be}\label{what-an-arbitrary-remaining-ale-core-would-have-to-be}}

There is also a classification-free part of this argument. If a
connected compact oriented four-manifold \(B\) has finite fundamental
group and connected rational-homology-three-sphere boundary and embeds
in a manifold satisfying (H), then \[
                         H_2(B;\mathbb Q)=0.             \tag{2a}
\] Indeed the map \(H_2(B;\mathbb Q)\to H_2(B,\partial B;\mathbb Q)\) is
an isomorphism by the boundary homology sequence. Poincaré--Lefschetz
duality therefore makes the rational intersection pairing nondegenerate.
The lifted embedding into \(\widetilde M\), whose second homology
vanishes, makes every such pairing zero. Nondegeneracy implies (2a).
Also \(H_1(B;\mathbb Q)=0\) by finiteness of its fundamental group.
Connected boundary and duality give \(H_3(B;\mathbb Q)=0\), so \(B\) is
a rational homology ball.

In particular, a remaining smooth Ricci-flat ALE bubble (these have
finite fundamental group and spherical-space-form ends) could not carry
nonzero rational second homology. Combining this with Theorem 1 and the
\href{20260923_critical_chiral_threshold.md}{paired-concentration
theorem}, a surviving bubble would have to be a \textbf{rational
homology ball with both Weyl halves nonzero and carrying positive
critical mass}. This does not claim that such Ricci-flat ALE spaces
exist, or exclude them by importing the classification restricted to
half-flat spaces. It narrows precisely the unresolved compact-regime
degeneration.

\hypertarget{the-enlarged-class-contains-genuinely-nonaspherical-manifolds}{%
\subsection{2. The enlarged class contains genuinely nonaspherical
manifolds}\label{the-enlarged-class-contains-genuinely-nonaspherical-manifolds}}

\textbf{Proposition 2.} Suppose \(X_1,\ldots,X_N\) are closed oriented
four-manifolds satisfying (H). Then \[
 M=X_1\#\cdots\#X_N\#m(S^1\times S^3)
 \tag{3}
\] also satisfies (H).

\textbf{Proof.} Van Kampen identifies \(\pi_1(M)\) with the free product
of the \(\pi_1(X_i)\) and a free group of rank \(m\). A finite-order
element of a free product is conjugate into one of its factors; this
group is therefore torsion-free.

Cut (3) along its connected-sum three-spheres. Its universal cover is a
tree of universal covers of punctured factors, glued along
three-spheres. Removing a locally finite collection of open four-balls
does not change \(H_1\) or \(H_2\): excision gives relative groups only
in degree four. Consequently each lifted vertex piece has vanishing
\(H_1,H_2\). The same holds for the punctured universal cover of
\(S^1\times S^3\), whose unpunctured cover is \(\mathbb R\times S^3\).
Mayer--Vietoris over finite subtrees, followed by the direct limit,
gives \(H_2(\widetilde M;\mathbb Z)=0\); edge three-spheres have no
first or second homology. Compact singular cycles involve only finitely
many pieces. Hurewicz now gives \(\pi_2(M)=0\). \(\square\)

For a concrete distinction, take \[
 M=(\Sigma_g\times\Sigma_h)\#m(S^1\times S^3),
 \qquad g,h\ge2,\quad m>0.                            \tag{4}
\] It satisfies (H), but is not aspherical. If it were aspherical it
would be a \(K(\pi,1)\) for \(\pi=\pi_1(\Sigma_g\times\Sigma_h)*F_m\).
Another finite \(K(\pi,1)\) is
\((\Sigma_g\times\Sigma_h)\vee\bigvee_{i=1}^m S^1\), of Euler
characteristic \(\chi(\Sigma_g\times\Sigma_h)-m\). In contrast the
connected sum (4) has Euler characteristic
\(\chi(\Sigma_g\times\Sigma_h)-2m\). Homotopy invariance contradicts
\(m>0\). Thus Theorem 1 is a genuine enlargement of the topological
scope, not another definition of asphericity.

\hypertarget{consequences-for-compactness-and-exact-realization}{%
\subsection{3. Consequences for compactness and exact
realization}\label{consequences-for-compactness-and-exact-realization}}

\textbf{Corollary 3.} Every assertion in
\href{20260923_compact_realization.md}{compact-regime realization},
including its critical equiintegrability version, remains valid with
``aspherical'' replaced by (H), with the same stated geometric
dependencies.

\textbf{Proof.} In that proof the geometric inputs are Ricci bounds,
diameter, volume, and one-sided Weyl control. Harmonic-radius blow-up
produces a smooth nonflat half-flat Ricci-flat ALE space and embeds its
compact core in the original manifold. Asphericity enters only in
forbidding that embedding. Apply Theorem 1 instead. The later
diffeomorphism compactness and Einstein-limit argument are unchanged.
\(\square\)

In particular, for \(q>2\), fixed \(D,v,\Lambda>0\), and topology (H),
\[
 \operatorname{diam}(g)\le D,\quad V_g\ge v,\quad
 \|W_g^+\|_{L^q}\le\Lambda,\quad
 \|\operatorname{Ric}^{\sharp}_g+I\|_{\infty,\mathrm{op}}
       <\varepsilon(D,v,\Lambda,q)
 \ \Longrightarrow\ \exists h:\operatorname{Ric}_h=-h.
 \tag{5}
\] The conclusion concerns the same smooth manifold. It does not assume
that the input is close to a previously known Einstein metric.

The theorem also supplies the topological step in the new
\href{20260923_critical_chiral_threshold.md}{critical chiral threshold}
and \href{20260923_gauge_saturation_ke.md}{gauge-saturation
realization}. Those notes prove their analytical inputs independently of
Proposition 2.

\hypertarget{an-application-to-handle-additions-with-actual-gauge-supply}{%
\subsection{4. An application to handle additions with actual gauge
supply}\label{an-application-to-handle-additions-with-actual-gauge-supply}}

For (4), let \(A=8(g-1)(h-1)>0\). The canonical SW class on the product,
followed by the handle surgery formula retaining the degree-one
insertions, supplies an actual conformally solvable class with \[
 a^2=A,\qquad (a_g^+)^2\ge A.
 \tag{6}
\] The supply statement, not merely this intersection calculation, is
the classical input detailed in the
\href{../mathematics/current/lebrun_stability_closure.md}{polarized
closure}, Sections 1--2. Here \(b_+>1\), so no chamber choice is
necessary.

If \(m\) is odd, then \(b_1(M)=2g+2h+m\) is odd. No complex structure on
this smooth manifold can be Kähler, in either orientation, since a
compact Kähler manifold has even first Betti number. Apply the
gauge-saturation realization theorem with its extension (H). For every
fixed diameter bound \(D\), it follows that \[
 \delta_-(g)+\frac{\mathfrak D_2(g,a)}{V_g}
       \ge\varepsilon(D,A_0)>0,
 \qquad (a_g^+)^2\ge A_0>0,\quad \operatorname{diam}(g)\le D.
 \tag{7}
\] One may choose \(A_0=8\) simultaneously for all products in (4). When
the sum on the left is small, the scalar monopole bound supplies the
volume lower bound used by the realization theorem; it is not an
additional hypothesis on these examples.

For instance \(g=2,h=11,m=1\) gives \(A=80\), \((2\chi+3\tau)(M)=76\),
and handle ratio \(4m/A=1/20\). The old displayed one-third-band
inequality with constant at most \(1/4\) is vacuous at zero defect for
this ratio, since \(1/3-1/20=17/60>1/4\). Statement (7) is a different
conclusion: it forbids joint gauge saturation and pointwise Einstein
approximation under a diameter bound. It does not replace that old
unrestricted integral gap or prove a gap for \(\delta_-\) alone.

No assertion is made that every topology in this infinite family admits
metrics satisfying a common diameter bound and very small Ricci defect.
Indeed (6) and volume comparison restrict the available topologies when
\(D\) is fixed. Nor does (7) exclude all non-Kähler Einstein metrics:
such a metric could have positive \(\mathfrak D_2/V\). Its precise
increment is an application of a new exact-realization implication to
nonaspherical smooth topology.

\hypertarget{dependencies-and-further-scope}{%
\subsection{5. Dependencies and further
scope}\label{dependencies-and-further-scope}}

Theorem 1 uses \href{https://arxiv.org/abs/1102.2442}{Wright,
\emph{Quotients of gravitational instantons}} and
\href{https://arxiv.org/abs/1102.2672}{Suvaina, \emph{ALE Ricci-flat
Kähler metrics and deformations of quotient surface singularities}}. The
actual gauge supplier in Section 4 is the one in
\href{https://arxiv.org/html/math/0003068}{LeBrun, Proposition 3.2},
with the handle insertions recorded in the inherited closure proof.

The condition (H) is sufficient, not proved necessary. More generally
the same argument works whenever finite core groups map trivially into
\(\pi_1(M)\) and the lifted ambient has the two homology vanishings used
above. Manifolds with nonzero \(\pi_2\), torsion in \(\pi_1\), or
unbounded-diameter collapse are not settled by this extension.

\end{document}
